\chapter{Geometria algebraiczna a analiza zespolona: analityfikacja, Chow, GAGA i Kähler}
\label{ch:ga9-analiza-zespolona}

\section{Dwie geometrie tych samych równań}
\label{sec:ga9-motywacja}

Przez kilka poprzednich rozdziałów badaliśmy przestrzenie opisane równaniami
wielomianowymi. Gdy ciało bazowe jest równe $\C$, możemy tę samą przestrzeń
oglądać za pomocą funkcji holomorficznych. Wówczas pojawia się pytanie:
\emph{jak wiele informacji traci się, zastępując wielomiany funkcjami
holomorficznymi?} Odpowiedź na przestrzeni afinicznej jest zupełnie inna niż
na przestrzeni rzutowej. Ta różnica tłumaczy sens twierdzeń Chowa i GAGA.

W całym rozdziale wszystkie schematy algebraiczne są skończonego typu nad
$\C$. Gdy mówimy o rozmaitości algebraicznej bez dodatkowego określenia,
zakładamy, że jest zredukowana i całkowita. W twierdzeniach dotyczących
rozmaitości zespolonych wskazujemy osobno warunki gładkości oraz zwartości.
Nie utożsamiamy przestrzeni topologicznej z jej snopem strukturalnym.

\begin{przyklad}[Pierścień funkcji regularnych i holomorficznych]
Na afinicznej prostej algebraicznej $\A^1_\C$ funkcje regularne na całej
przestrzeni są wielomianami: $\mathcal O_{\rm alg}(\A^1)=\C[z]$.
Na zespolonej płaszczyźnie $\C$ funkcje holomorficzne tworzą znacznie
większy pierścień $\mathcal O_{\rm hol}(\C)$; należy do niego $e^z$.
Nie istnieje wielomian $P\in\C[z]$ taki, że $P(z)=e^z$ dla każdego
$z\in\C$. Istotnie, po ograniczeniu do dodatniej półosi rzeczywistej
otrzymalibyśmy $P(x)=e^x$, lecz $e^x/x^m\to+\infty$ dla każdego
$m\ge0$, podczas gdy $|P(x)|/x^m$ pozostaje ograniczone, jeśli
$m\ge\deg P$. Jest to przykład różnicy między funkcjami regularnymi
i holomorficznymi, a nie tylko różnicy między wybranymi wzorami.
\end{przyklad}

\begin{uwaga}[Ten sam zbiór punktów nie oznacza tej samej kategorii]
Punkty zespolone $\A^1(\C)$ i punkty analityczne $\C$ można utożsamić,
lecz w topologii Zariskiego zbiory otwarte są dopełnieniami zbiorów
skończonych, podczas gdy w topologii euklidesowej otwartych zbiorów jest
znacznie więcej. Co ważniejsze, różnią się też odpowiednie snopy
funkcji. Analityfikacja jest więc konstrukcją na przestrzeniach
\emph{z ich strukturą}, nie tylko zmianą nazwy punktów.
\end{uwaga}

\section{Zbiory analityczne i analityfikacja odmian}
\label{sec:ga9-analityczne}

\begin{definicja}[Zbiór analityczny]
Niech $M$ będzie rozmaitością zespoloną o wymiarze zespolonym $n$,
a $A\subseteq M$ podzbiorem. Mówimy, że $A$ jest
\emph{zespolonym zbiorem analitycznym}, jeżeli dla każdego $p\in M$
istnieje otoczenie otwarte $U$ punktu $p$ i funkcje holomorficzne
$f_1,\ldots,f_r\in\mathcal O_M(U)$ takie, że
$A\cap U=\{x\in U:f_1(x)=\cdots=f_r(x)=0\}$.
Jeśli lokalnie rozpatrujemy dodatkowo snop ideałów funkcji znikających,
otrzymujemy odpowiednią zredukowaną przestrzeń analityczną.
\end{definicja}

\begin{uwaga}[Lokalność nie oznacza globalnego układu równań]
Definicja dopuszcza różne równania na różnych otoczeniach $U$.
W ogólności zbiór analityczny w rozmaitości zespolonej nie musi być
zbiorem wspólnych zer skończonej rodziny globalnych funkcji
holomorficznych. Na przykład na zwartej spójnej rozmaitości zespolonej
wszystkie globalne funkcje holomorficzne są stałe, lecz mogą istnieć
nietrywialne dzielniki analityczne. Snopowość opisu jest zatem
niezbędna.
\end{uwaga}

\begin{definicja}[Analityfikacja]
Niech $X$ będzie schematem skończonego typu nad $\C$. Jego
\emph{analityfikacją} nazywamy zespoloną przestrzeń analityczną
$X^{\rm an}$, której punkty odpowiadają punktom $X(\C)$,
wraz z kanonicznym odwzorowaniem $a:X^{\rm an}\to X$
do przestrzeni Zariskiego i morfizmem snopów
$a^{-1}\mathcal O_X\to\mathcal O_{X^{\rm an}}$.
Jej lokalny model dla
$X=\operatorname{Spec}\C[z_1,\ldots,z_n]/I$ otrzymujemy jako wspólny
zbiór zer generatorów $I$ w $\C^n$, wyposażony w iloraz snopu
funkcji holomorficznych przez snop ideałów generowany przez $I$.
Modele afiniczne skleja się przy użyciu holomorficznych przejść
pochodzących od morfizmów algebraicznych. Dla gładkiego $X$ jest to
zwykła rozmaitość zespolona.
\end{definicja}

\begin{uwaga}[Nilpotenty i osobliwości]
Jeżeli $I=(w^2)\subset\C[z,w]$, to analityfikacja schematu
$\operatorname{Spec}\C[z,w]/(w^2)$ ma ten sam zbiór punktów co prosta
$w=0$, lecz jej snop strukturalny nadal zawiera niezerowy nilpotent
$w$. Analityfikacja schematów nie polega więc na wyrzucaniu informacji
o krotnościach. W przypadku $I=(w^2-z^3)$ otrzymujemy ten sam
analityczny kiełek ostrza, którego normalizację opisaliśmy w
rozdziałach~\ref{ch:ga5-rozdmuchania} i~\ref{ch:ga6-powierzchnie}.
\end{uwaga}

\begin{stwierdzenie}[Wielomiany definiują zbiory analityczne]
Niech $I\subseteq\C[z_1,\ldots,z_n]$ będzie ideałem. Wówczas
$V(I)\subseteq\C^n$, rozpatrywany w zwykłej topologii, jest
zamkniętym zbiorem analitycznym. Jeśli $I$ jest generowany przez $f_1,\ldots,f_r$, to te same
wielomiany definiują analitycznie \emph{zbiór punktów} $V(I)$.
Aby otrzymać analityfikację schematu z ewentualnymi nilpotentami,
należy ponadto zachować iloraz snopu holomorficznego przez snop
ideałów generowany przez $f_1,\ldots,f_r$.
\end{stwierdzenie}
\begin{proof}
Każdy wielomian jest funkcją holomorficzną, a wspólny zbiór zer
funkcji ciągłych jest domknięty. Ponieważ $\C[z_1,\ldots,z_n]$ jest
pierścieniem noetherowskim (twierdzenie Hilberta o bazie z
rozdziału~\ref{ch:pierscienie}), ideał $I$ ma skończony układ
generatorów. Definicję zbioru analitycznego spełnia więc nawet
jedno otoczenie $U=\C^n$ z tym właśnie układem równań.
\end{proof}

\section{Dlaczego w przestrzeni afinicznej odwrotność zawodzi?}
\label{sec:ga9-wykladnicza}

\begin{przyklad}[Analityczna, lecz niealgebraiczna krzywa]
\label{prz:ga9-exp}
Rozważmy $\Gamma=\{(z,w)\in\C^2:w=e^z\}$.
Jest to zamknięty zbiór analityczny, ponieważ $w-e^z$ jest
funkcją holomorficzną, a różniczka $\mathrm d(w-e^z)$ ma współczynnik
przy $\mathrm dw$ równy $1$; zatem $\Gamma$ jest nawet gładką
podrozmaitością zespoloną wymiaru jeden. Pokażemy, że nie jest
zbiorem algebraicznym.

Załóżmy, że niezerowy wielomian $P(z,w)=\sum_{j=0}^m p_j(z)w^j$
znika na $\Gamma$, gdzie $p_m\ne0$. Wtedy dla każdego $z\in\C$
mamy $\sum_{j=0}^m p_j(z)e^{jz}=0$. Wstawmy $z=x\in\R$
i podzielmy przez $e^{mx}$. Otrzymujemy
\[
 p_m(x)=-\sum_{j<m}p_j(x)e^{-(m-j)x}.
\]
Prawa strona dąży do zera, gdy $x\to+\infty$, gdyż każda
funkcja wykładnicza $e^{-cx}$, $c>0$, dominuje wielomianowy
wzrost $|p_j(x)|$. Oznaczałoby to $p_m(x)\to0$, co jest niemożliwe
dla niezerowego wielomianu. Zatem żaden niezerowy wielomian
nie znika na całej $\Gamma$. W szczególności $\Gamma$ nie może
być algebraiczną krzywą w $\A^2_\C$.
\end{przyklad}

\begin{center}
\begin{tikzpicture}[>=Stealth,scale=.93]
  \draw[->] (-3.2,0)--(3.15,0) node[right,font=\footnotesize] {$x$};
  \draw[->] (0,-.35)--(0,3.4) node[above,font=\footnotesize] {$w$};
  \draw[blue!75!black,very thick,domain=-3.1:1.18,samples=80]
    plot (\x,{.65*exp(\x)});
  \node[fill=white,inner sep=2pt,font=\footnotesize,anchor=west]
    at (.72,2.98) {$w=e^x$};
  \node[font=\scriptsize,align=center] at (-1.65,2.65)
    {rzeczywisty przekrój\\zespolonej krzywej $\Gamma$};
\end{tikzpicture}

{\small Rysunek 9.1. Rzeczywisty ślad grafu funkcji wykładniczej;
rysunek jest wyłącznie przekrojem pomocniczym, a nie modelem
pełnej krzywej zespolonej w $\C^2$.}\par
\end{center}

\begin{uwaga}[Dlaczego twierdzenie Chowa nie daje sprzeczności?]
Graf funkcji wykładniczej jest zamknięty w afinicznej przestrzeni
$\C^2$, ale jego obraz w $\PP^2(\C)$ nie jest domkniętym zbiorem
analitycznym. Twierdzenie Chowa będzie dotyczyło tylko zbiorów
\emph{zamkniętych analitycznie w przestrzeni rzutowej}.
Nie wystarcza sama algebraiczność przestrzeni otaczającej.
\end{uwaga}

\section{Lokalne pierścienie holomorficzne i przygotowanie Weierstrassa}
\label{sec:ga9-lokalne}

\begin{definicja}[Pierścień kiełków holomorficznych]
Niech $p\in\C^n$. Przez $\mathcal O_{\C^n,p}$ oznaczamy pierścień
klas par $(U,f)$, gdzie $p\in U\subseteq\C^n$ jest otwarty,
$f\in\mathcal O(U)$, a dwie pary są równoważne, gdy funkcje
zgadzają się na pewnym otoczeniu $p$. Jest to pierścień lokalny
o ideale maksymalnym $\mathfrak m_p=\{[f]:f(p)=0\}$.
Po przesunięciu $p$ do zera identyfikujemy go z pierścieniem
zbieżnych szeregów potęgowych $\C\{z_1,\ldots,z_n\}$.
\end{definicja}

\begin{proof}
Klasy dodajemy i mnożymy przez wybór wspólnego, odpowiednio
zmniejszonego otoczenia. Jeśli $f(p)\ne0$, ciągłość daje
otoczenie, na którym $f$ nie ma zer, więc $1/f$ jest tam
holomorficzna i $[f]$ jest jednostką. Jeśli $f(p)=0$,
to $[f]$ nie może być jednostką: wartość w $p$ iloczynu
$[f][g]$ jest równa zero. Elementy nieodwracalne tworzą
zatem dokładnie ideał $\mathfrak m_p$, który jest maksymalny,
bo homomorfizm wartościowania $[f]\mapsto f(p)$ ma jako
jądro $\mathfrak m_p$ i jako obraz $\C$. Jedyność ideału
maksymalnego wynika stąd, że każdy właściwy ideał składa się
z elementów nieodwracalnych. Identyfikacja z pierścieniem
szeregów zbieżnych wynika z lokalnego rozwinięcia Taylora
funkcji holomorficznej.
\end{proof}

\begin{twierdzenie}[Przygotowanie Weierstrassa — wynik zewnętrzny]
\label{tw:ga9-weierstrass}
Niech $n\ge1$ będzie liczbą całkowitą,
$f\in\C\{z_1,\ldots,z_n\}$ i załóżmy, że
$f(0,\ldots,0,z_n)$ ma w zerze rząd dokładnie $d\ge1$.
Wtedy istnieje jednostka $u\in\C\{z_1,\ldots,z_n\}^\times$
oraz jednoznaczny wielomian wyróżniony
\[
 W=z_n^d+a_1(z')z_n^{d-1}+\cdots+a_d(z'),
 \qquad a_j\in\C\{z'\},\quad a_j(0)=0,
\]
gdzie $z'=(z_1,\ldots,z_{n-1})$, taki że $f=uW$.
\end{twierdzenie}

\begin{uwaga}[Znaczenie twierdzenia]
Choć $f$ może mieć nieskończenie wiele niezerowych wyrazów
szeregu Taylora, po pomnożeniu przez jednostkę można traktować
jego zera lokalnie jak zera \emph{wielomianu w jednej zmiennej}
o współczynnikach holomorficznych w pozostałych zmiennych.
Jest to punkt wyjścia do finitości i lokalnej teorii zbiorów
analitycznych. Dowód przygotowania Weierstrassa jest ważnym,
samodzielnym rezultatem analizy zespolonej; tutaj podajemy jego
precyzyjną treść, nie zastępując dowodu nieuzasadnionym
rachunkiem formalnym.
\end{uwaga}

\begin{przyklad}[Ostrze w pierścieniu zbieżnych szeregów]
Funkcja $f(z,w)=w^2-z^3$ jest już wielomianem wyróżnionym
stopnia dwa względem $w$. Stąd lokalny pierścień analityczny
ostrza w zerze jest $\C\{z,w\}/(w^2-z^3)$.
Parametryzacja $z=t^2$, $w=t^3$ daje naturalny homomorfizm
$\C\{z,w\}/(w^2-z^3)\to\C\{t\}$, którego obrazem jest
$\C\{t^2,t^3\}$. Po dołączeniu funkcji $t$ otrzymujemy
normalizację w sensie analitycznym. Ta sama parametryzacja
działa algebraicznie, lecz pierścienie $\C[t]$ i $\C\{t\}$
nie są tym samym obiektem.
\end{przyklad}

\section{Snopy koherentne: właściwy język porównania}
\label{sec:ga9-koherentne}

\begin{definicja}[Koherentny snop analityczny]
Niech $X$ będzie zespoloną przestrzenią analityczną,
a $\mathcal F$ snopem modułów nad $\mathcal O_X$.
Nazywamy $\mathcal F$ \emph{koherentnym}, jeżeli lokalnie
jest skończenie prezentowany i jądra morfizmów z lokalnie
wolnych snopów skończonego rzędu do $\mathcal F$ są lokalnie
skończenie generowane. Równoważnie, w standardowej teorii
snopów koherentnych nad zespolonymi przestrzeniami
analitycznymi, $\mathcal F$ jest lokalnie ilorazem
$\mathcal O_X^{\oplus r}$ przez koherentny snop relacji.
\end{definicja}

\begin{uwaga}[Dlaczego potrzebny jest warunek na relacje?]
Przy przejściu do obrazu lub jądra morfizmu snopów trzeba
kontrolować nie tylko liczbę generatorów, ale również ich
relacje. Fundamentalne twierdzenie Oki o koherencji snopu
$\mathcal O_{\C^n}$ oraz jego konsekwencje zapewniają
stabilność koherentności przy operacjach występujących
w geometrii analitycznej. Są to wyniki zewnętrzne.
\end{uwaga}

\begin{przyklad}[Snop ideałów podprzestrzeni analitycznej]
Jeśli $A\subseteq X$ jest zamkniętym zredukowanym
podzbiorem analitycznym, przez $\mathcal I_A$ oznaczamy
snop kiełków funkcji holomorficznych zerujących się na $A$.
Wtedy $\mathcal O_A=\mathcal O_X/\mathcal I_A$,
a ciąg
$0\to\mathcal I_A\to\mathcal O_X\to\mathcal O_A\to0$
jest dokładny. Koherentność $\mathcal I_A$ jest
konsekwencją twierdzeń Oki--Cartana; nie wynika wyłącznie
z definicji analityczności przez lokalne równania.
\end{przyklad}

\section{Twierdzenie GAGA: co dokładnie porównujemy?}
\label{sec:ga9-gaga}

\begin{twierdzenie}[Serre, GAGA — wynik zewnętrzny]
\label{tw:ga9-gaga}
Niech $X$ będzie schematem właściwym i skończonego typu nad $\C$.
Analityfikacja wyznacza równoważność kategorii
\[
 \operatorname{Coh}(X)\ \simeq\ \operatorname{Coh}(X^{\rm an}),
 \qquad \mathcal F\longmapsto\mathcal F^{\rm an},
\]
między koherentnymi snopami algebraicznymi i koherentnymi
snopami analitycznymi. Ponadto dla każdego koherentnego
snopu algebraicznego $\mathcal F$ oraz $q\ge0$ naturalny
homomorfizm
$H^q(X,\mathcal F)\to H^q(X^{\rm an},\mathcal F^{\rm an})$
jest izomorfizmem. W szczególności każdy koherentny
snop analityczny na $X^{\rm an}$ ma, z dokładnością do
izomorfizmu, jednoznaczny algebraiczny odpowiednik.
\end{twierdzenie}

\begin{uwaga}[Co znaczy równoważność, a czego ona nie mówi?]
Twierdzenie obejmuje nie tylko istnienie algebraicznego
snopu, lecz także pełną wierność: wszystkie morfizmy
koherentnych snopów analitycznych między analityfikacjami
pochodzą dokładnie z morfizmów algebraicznych.
Nie oznacza to równości pierścieni funkcji holomorficznych
i regularnych na dowolnym \emph{afinicznym} otwartym
podzbiorze $X^{\rm an}$: warunek właściwości jest
istotny. Nie oznacza też, że bezpośrednio porównujemy
wszystkie snopy zbiorów, tylko snopy koherentne.
\end{uwaga}

\begin{center}
\begin{tikzpicture}[>=Stealth,
 every node/.style={font=\footnotesize,align=center},
 nodebox/.style={draw,rounded corners,inner sep=6pt,text width=46mm}]
 \node[nodebox] (a) at (-3.2,1.3)
   {snopy koherentne\\algebraiczne na $X$};
 \node[nodebox] (b) at (3.2,1.3)
   {snopy koherentne\\analityczne na $X^{\rm an}$};
 \node[nodebox] (c) at (-3.2,-1.05)
   {$H^q(X,\mathcal F)$};
 \node[nodebox] (d) at (3.2,-1.05)
   {$H^q(X^{\rm an},\mathcal F^{\rm an})$};
 \draw[<->,thick] (a)--node[above,font=\scriptsize] {GAGA} (b);
 \draw[->] (a)--node[left,font=\scriptsize] {kohomologia} (c);
 \draw[->] (b)--node[right,font=\scriptsize] {kohomologia} (d);
 \draw[<->,thick] (c)--(d);
\end{tikzpicture}

{\small Rysunek 9.2. Dwa porównania twierdzenia GAGA.
Górna strzałka oznacza równoważność kategorii, dolna
izomorfizm grup; nie są to identyfikacje definicji.}\par
\end{center}

\begin{stwierdzenie}[Porównanie wiązek liniowych]
Niech $X$ będzie schematem właściwym nad $\C$.
Analityfikacja wyznacza izomorfizm grup
$\operatorname{Pic}(X)\simeq\operatorname{Pic}(X^{\rm an})$.
W szczególności na zespolonej rozmaitości rzutowej
każda holomorficzna wiązka liniowa pochodzi z algebraicznej.
\end{stwierdzenie}
\begin{proof}
Wiązka liniowa to koherentny snop lokalnie wolny rzędu jeden.
GAGA daje algebraiczny snop koherentny $\mathcal L$ odpowiadający
analitycznej wiązce liniowej. Lokalna wolność $\mathcal L$
może być testowana po analityfikacji, ponieważ homomorfizmy
lokalnych pierścieni algebraicznych do holomorficznych
pierścieni lokalnych są wiernie płaskie; jest to standardowa
lokalna część twierdzenia porównawczego. Stąd $\mathcal L$
jest algebraiczną wiązką liniową. Iloczyn tensorowy
i wiązka odwrotna są zachowywane przez analityfikację,
co daje izomorfizm grup Picarda.
\end{proof}

\begin{uwaga}[Gdzie ukrywa się trudność dowodu GAGA?]
Wśród głębokich składników pełnego dowodu znajdują się
analityczna skończoność kohomologii snopów koherentnych
na przestrzeniach zwartych oraz algebraiczne twierdzenia
Serre'a o snopach koherentnych na przestrzeniach rzutowych.
W tym skrypcie udowodnimy ich najważniejsze konsekwencje
na $\PP^1$ bez odwoływania się do pełnej teorii; samą
równoważność GAGA przyjmujemy jako twierdzenie zewnętrzne.
\end{uwaga}

\section{Chow: dlaczego zamknięta geometria analityczna staje się algebraiczna?}
\label{sec:ga9-chow}

\begin{twierdzenie}[Chow]
\label{tw:ga9-chow}
Niech $n\ge0$ będzie liczbą całkowitą, a
$A\subseteq\PP^n(\C)$ podzbiorem zamkniętym
w zwykłej topologii zespolonej, który
jest zespolonym zbiorem analitycznym. Wtedy istnieją
jednorodne wielomiany $F_1,\ldots,F_r\in\C[Z_0,\ldots,Z_n]$
takie, że
$A=\{[Z]\in\PP^n(\C):F_1(Z)=\cdots=F_r(Z)=0\}$.
Zatem $A$ jest rzutowym zbiorem algebraicznym.
\end{twierdzenie}
\begin{proof}[Dowód jako konsekwencja GAGA]
Wprowadzimy jawnie wyniki użyte w przejściu.
Zbiór $A$ ma koherentny analityczny snop ideałów
$\mathcal I_A\subseteq\mathcal O_{\PP^n{}^{\rm an}}$
(twierdzenie Oki--Cartana). Z twierdzenia GAGA
\ref{tw:ga9-gaga} wynika istnienie koherentnego
algebraicznego snopu $\mathcal I$ na $\PP^n_\C$
oraz izomorfizmu $\mathcal I^{\rm an}\simeq\mathcal I_A$.
Ponieważ analityfikacja jest w GAGA w pełni wierna,
inkluzja $\mathcal I_A\hookrightarrow\mathcal O_{\PP^n{}^{\rm an}}$
pochodzi z morfizmu algebraicznego
$\mathcal I\to\mathcal O_{\PP^n}$. Jest to iniekcja:
można to sprawdzić po wiernie płaskim przejściu do
lokalnych pierścieni analitycznych. Tak więc
$\mathcal I$ jest koherentnym snopem ideałów.
Wyznacza on zamknięty podschemat
$Z\hookrightarrow\PP^n_\C$.

Pozostaje wykazać, że algebraiczny podschemat rzutowy
jest opisany przez skończenie wiele jednorodnych
wielomianów. Na mocy twierdzeń Serre'a o snopach
na przestrzeni rzutowej dla dostatecznie dużego $d$
snop $\mathcal I(d)$ jest generowany przez skończenie
wiele sekcji globalnych. Każdą z nich można
interpretować jako jednorodny wielomian stopnia $d$;
gdy potrzebne są różne stopnie, wybieramy skończony
zestaw takich równań we wszystkich mapach współrzędnych
(z noetherowskości i koherencji). Powstający
jednorodny ideał ma jako związany snop ideałów
$\mathcal I$. W szczególności jego zbiór punktów
zespolonych jest dokładnie $A$. To kończy dowód.
\end{proof}

\begin{uwaga}[Status dowodu]
Dowód jest pełnym \emph{wyprowadzeniem Chow z GAGA,
koherencji analitycznej i twierdzeń Serre'a o przestrzeni
rzutowej}. Nie jest elementarnym dowodem Chowa od podstaw:
wymienione trzy rezultaty przyjmujemy jako wyniki zewnętrzne.
Takie rozgraniczenie jest ważne, by nie ukrywać trudności
we frazie „z teorii snopów wynika”.
\end{uwaga}

\begin{wniosek}[Holomorficzne odwzorowania projektowe są algebraiczne]
\label{wn:ga9-hol-map}
Niech $X,Y$ będą zespolonymi rzutowymi rozmaitościami
algebraicznymi (niekoniecznie gładkimi). Każde
odwzorowanie holomorficzne
$f:X^{\rm an}\to Y^{\rm an}$ jest analityfikacją
jednego morfizmu algebraicznego $X\to Y$.
\end{wniosek}
\begin{proof}
Graf $\Gamma_f\subseteq X^{\rm an}\times Y^{\rm an}$
jest zamkniętym zbiorem analitycznym, ponieważ
$Y^{\rm an}$ jest rozdzielona, a lokalnie graf
definiują równania holomorficzne odwzorowania. Ponieważ
$X\times Y$ jest rzutowy (zanurzenie Segrego),
twierdzenie Chowa pokazuje, że $\Gamma_f$ jest
algebraicznym podzbiorem $X\times Y$. Projekcja
$\Gamma_f\to X$ jest izomorfizmem po analityfikacji.
Algebraiczny morfizm właściwy, którego analityfikacja
jest izomorfizmem, jest izomorfizmem (konsekwencja
lokalnej wiernej płaskości i GAGA). Dlatego
projekcja algebraiczna $\Gamma_f\to X$ jest
izomorfizmem. Złożenie jej odwrotności z drugą
projekcją do $Y$ daje poszukiwany morfizm.
Jednoznaczność wynika z pełnej wierności porównania morfizmów
w sytuacji właściwej (lub, dla zredukowanych odmian, z faktu,
że morfizmy do rozdzielonego celu są wyznaczone przez swoje
wartości na gęstym zbiorze punktów domkniętych).
\end{proof}

\begin{przyklad}[Dlaczego nie wolno pominąć rzutowości?]
Holomorficzne odwzorowanie $\C\to\C^*$, $z\mapsto e^z$,
nie jest morfizmem algebraicznym. Morfizm algebraiczny
$\A^1_\C\to\mathbb G_m$ odpowiada jednostce w $\C[z]$,
a jedynymi jednostkami tego pierścienia są niezerowe
stałe. Zatem w tym przypadku wszystkie morfizmy
algebraiczne do $\mathbb G_m$ są stałe, mimo istnienia wielu
niestacjonarnych odwzorowań holomorficznych.
\end{przyklad}

\section[Rachunek porównawczy na P1: kohomologia w obu językach]{Rachunek porównawczy na $\PP^1$: kohomologia w obu językach}
\label{sec:ga9-p1-cech}

Rozpatrzmy standardowe pokrycie $\PP^1=U_0\cup U_1$,
gdzie $U_0\cong\A^1$ ma współrzędną $z=Z_1/Z_0$,
a $U_1\cong\A^1$ ma współrzędną $w=Z_0/Z_1=1/z$.
Na przecięciu dziedzin obu map współrzędnych $z\ne0$.
Niech $e_0,e_1$ będą standardowymi lokalnymi bazami
wiązki $\mathcal O(n)$, przy konwencji $e_1=z^n e_0$.

\begin{stwierdzenie}[Kohomologia algebraiczna $\mathcal O(n)$]
\label{st:ga9-p1-coh}
Niech $n\in\Z$. Dla wiązki liniowej $\mathcal O_{\PP^1_\C}(n)$
na prostej rzutowej nad $\C$ mamy
\[
 h^0(\PP^1,\mathcal O(n))=\max\{n+1,0\},\qquad
 h^1(\PP^1,\mathcal O(n))=\max\{-n-1,0\}.
\]
Ponadto $H^q(\PP^1,\mathcal O(n))=0$ dla $q\ge2$.
\end{stwierdzenie}
\begin{proof}
Część algebraiczna kompleksu Čecha dla pokrycia
$U_0,U_1$ ma postać
\[
 \C[z]\oplus\C[z^{-1}]
 \ \xrightarrow{\ \delta\ }\ \C[z,z^{-1}],\qquad
 \delta(a,b)=a-z^n b.
\]
Znak nie ma wpływu na iloraz, ale określamy go,
aby ustalić konwencję trywializacji. Sekcja globalna
to para $(a,b)$ z $a=z^n b$. Jeśli $n\ge0$,
możemy przyjąć $b=\sum_{j=0}^n b_jz^{-j}$,
co daje wielomiany $a$ stopnia najwyżej $n$;
wymiar przestrzeni sekcji wynosi $n+1$. Gdy $n<0$,
nie ma niezerowej pary spełniającej równanie.

Grupa $H^1$ jest ilorazem
$\C[z,z^{-1}]/(\C[z]+z^n\C[z^{-1}])$.
W mianowniku znajdują się wszystkie potęgi $z^j$
dla $j\ge0$ oraz wszystkie potęgi $z^j$ dla $j\le n$.
Gdy $n\ge-1$, nie pozostaje żadna potęga. Gdy
$n\le-2$, bazę ilorazu tworzą klasy
$z^{n+1},z^{n+2},\ldots,z^{-1}$, których jest $-n-1$.
Ponieważ $U_0,U_1,U_0\cap U_1$ są afiniczne,
kohomologia koherentna pokrycia jest liczona
przez kompleks Čecha. Wymiar $\PP^1$ wynosi jeden,
więc dla snopów koherentnych wyższe grupy $H^q$
znikają; można to również odczytać z dwuczłonowego
kompleksu Čecha i acykliczności tego pokrycia.
\end{proof}

\begin{stwierdzenie}[Rachunek analityczny daje ten sam wynik]
\label{st:ga9-p1-analytic}
Niech $\mathcal O(n)^{\rm an}$ będzie analityfikacją
wiązki $\mathcal O(n)$ na sferze Riemanna
$\PP^1(\C)$. Grupy $H^0$ i $H^1$ tej wiązki
mają wymiary podane w
stwierdzeniu~\ref{st:ga9-p1-coh}.
\end{stwierdzenie}
\begin{proof}
W dziedzinach map analitycznych $U_0^{\rm an}=\C$ oraz
$U_1^{\rm an}=\C$ sekcje lokalne są już funkcjami
całkowitymi, a na przecięciu $\C^*$ -- funkcjami
holomorficznymi. Dla $H^0$ relacja
$a(z)=z^n b(1/z)$ wraz z rozwinięciami Taylora
w zerze i nieskończoności wymusza dokładnie
te same wielomiany co w dowodzie algebraicznym.

Każdą holomorficzną funkcję $g$ na $\C^*$ można
rozwinąć w szereg Laurenta zbieżny na
$0<|z|<\infty$. Część o wykładnikach nieujemnych
przedłuża się do funkcji całkowitej $a(z)$.
Część o wykładnikach $j\le n$ ma postać
$z^n b(1/z)$ z funkcją całkowitą $b$.
Pozostają tylko wykładniki $n<j<0$, jeśli takie
istnieją, a jest ich $-n-1$ dla $n\le-2$.
Zatem iloraz analitycznych $1$-kokykli przez
kobrzegi ma tę samą bazę co iloraz algebraiczny.
Aby utożsamić ten iloraz z analityczną grupą
$H^1$, wykorzystujemy acykliczność Steinowskich
otwartych zbiorów $U_0^{\rm an}$, $U_1^{\rm an}$
i ich przecięcia, które są Steinowskie (twierdzenie Cartana B).
Nie jest to niejawne zastosowanie samego GAGA.
\end{proof}

\begin{przyklad}[Dualność na sferze jako konkretny rachunek]
Dla $n=-3$ otrzymujemy $H^0(\PP^1,\mathcal O(-3))=0$
i $H^1(\PP^1,\mathcal O(-3))$ o bazie reprezentowanej
przez $z^{-2},z^{-1}$. Z drugiej strony
$H^0(\PP^1,\mathcal O(1))$ ma bazę $1,z$.
Parowanie residuowe z rozdziału~\ref{ch:krzywe-snopy}
łączy te przestrzenie w sposób niezdegenerowany;
odpowiada temu izomorfizm dualności Serre'a
$H^1(\mathcal O(-3))^*\cong H^0(\mathcal O(1))$.
Ten przykład pokazuje jednocześnie, jak wyniki GAGA
łączą algebraiczną teorię dzielników z klasycznymi
całkami po okręgach na sferze Riemanna.
\end{przyklad}

\section{Metryki hermitowskie i geometria Kählerowska}
\label{sec:ga9-kahler}

Na rozmaitości zespolonej $M$ mamy rozkład zespolonej
wiązki kostycznej na podwiązki generowane lokalnie przez
$\mathrm dz_1,\ldots,\mathrm dz_n$ i przez
$\mathrm d\bar z_1,\ldots,\mathrm d\bar z_n$.
Formę będącą kombinacją iloczynów zewnętrznych $p$ symboli
$\mathrm dz_j$ oraz $q$ symboli $\mathrm d\bar z_k$
nazywamy formą typu $(p,q)$. Zewnętrzna różniczka
$\mathrm d$ rozkłada się na $\partial$ i $\bar\partial$:
$\partial$ podnosi pierwszy stopień, a $\bar\partial$ drugi.
Na przykład $\bar\partial f=\sum_j
(\partial f/\partial\bar z_j)\,\mathrm d\bar z_j$;
funkcja klasy $C^1$ jest holomorficzna dokładnie wtedy,
gdy $\bar\partial f=0$. To wyjaśnia oznaczenia używane
w konstrukcji formy Fubiniego--Study'ego.

\begin{definicja}[Forma Kählerowska]
Niech $M$ będzie rozmaitością zespoloną wymiaru $n$
z operatorem struktury zespolonej $J:TM\to TM$.
Metryka Riemannowska $g$ jest \emph{hermitowska}, gdy
$g(Ju,Jv)=g(u,v)$ dla wszystkich rzeczywistych
wektorów stycznych $u,v$. Jej związana forma
$\omega(u,v)=g(Ju,v)$ jest rzeczywistą formą różniczkową
stopnia dwa typu $(1,1)$. Metrykę (lub formę)
nazywamy \emph{Kählerowską}, gdy dodatkowo
$\mathrm d\omega=0$. Rozmaitość $M$ nazywamy
\emph{Kählerowską}, jeśli dopuszcza taką metrykę.
\end{definicja}

\begin{przyklad}[Przestrzeń zespolona jako wzorzec]
Na $\C^n$ ze współrzędnymi $z_j=x_j+iy_j$
standardowa metryka euklidesowa ma formę
$\omega_0=\sum_{j=1}^n\mathrm dx_j\wedge\mathrm dy_j
=\frac{i}{2}\sum_{j=1}^n\mathrm dz_j\wedge\mathrm d\bar z_j$.
Jej współczynniki są stałe, więc $\mathrm d\omega_0=0$.
Każda kompleksowa podrozmaitość otwartego podzbioru
$\C^n$ dziedziczy przez restrykcję metrykę Kählerowską.
\end{przyklad}

\begin{twierdzenie}[Forma Fubiniego--Study'ego]
\label{tw:ga9-fs}
Niech $n\ge1$ będzie liczbą całkowitą i ustalmy standardowy iloczyn
hermitowski na $\C^{n+1}$. Na $\PP^n(\C)$ istnieje
wyznaczona przez niego dodatnia,
zamknięta rzeczywista $(1,1)$-forma
$\omega_{\rm FS}$, która na mapie $Z_0\ne0$
o współrzędnych $z_j=Z_j/Z_0$ jest dana wzorem
\[
 \omega_{\rm FS}=\frac{i}{2\pi}\partial\bar\partial
 \log(1+|z_1|^2+\cdots+|z_n|^2).
\]
Jej klasa de Rhama jest obrazem całkowitej klasy
$c_1(\mathcal O_{\PP^n}(1))$ w $H^2(\PP^n(\C),\R)$.
\end{twierdzenie}
\begin{proof}
Na dowolnym mapie rzutowym wybieramy reprezentanta
$Z=(Z_0,\ldots,Z_n)$ i stosujemy potencjał
$\log\sum_j|Z_j|^2$ po ustawieniu jednej współrzędnej
równej $1$. Dwie takie trywializacje różnią się
pomnożeniem reprezentanta przez niezerową funkcję
holomorficzną $h$, więc ich potencjały różnią się
o $\log|h|^2$. Ponieważ lokalnie
$\partial\bar\partial\log|h|^2=0$, otrzymane formy
sklejają się w globalną formę $\omega_{\rm FS}$.
Zamkniętość wynika z $\partial^2=\bar\partial^2=0$
i $\partial\bar\partial+\bar\partial\partial=0$.

Aby sprawdzić dodatniość, połóżmy $A=1+\|z\|^2$.
Macierz współczynników formy (z pominięciem dodatniego
czynnika $1/\pi$) ma postać
\[
 \frac{\partial^2\log A}{\partial z_j\partial\bar z_k}
 =\frac{A\delta_{jk}-\bar z_j z_k}{A^2}.
\]
Dla $v\ne0$ odpowiadająca jej forma hermitowska
jest proporcjonalna do
$\big(A\|v\|^2-|\sum_j\bar z_jv_j|^2\big)/A^2$.
Z nierówności Cauchy'ego--Schwarza dostajemy dolne
oszacowanie przez $\|v\|^2/A^2>0$.

Na końcu podajemy uzasadnienie identyfikacji klasy.
Tautologiczna wiązka $\mathcal O(-1)$ ma w lokalnej
bazie $(1,z_1,\ldots,z_n)$ metrykę hermitowską
o kwadracie normy $A$. Dualna wiązka $\mathcal O(1)$
ma kwadrat normy lokalnego dualnego przekroju równy $A^{-1}$.
Wzór Cherna--Weila
$c_1(L,h)=-(i/2\pi)\partial\bar\partial\log\|e\|^2_h$
dla lokalnej bazy $e$ daje dla $\mathcal O(1)$
właśnie $\omega_{\rm FS}$. Fakt, że forma krzywizny
reprezentuje topologiczną klasę Cherna, jest
standardowym twierdzeniem Cherna--Weila, przyjętym
tutaj jako wynik zewnętrzny. Wszystkie pozostałe
rachunki wykonaliśmy jawnie.
\end{proof}

\begin{przyklad}[Normalizacja całki na $\PP^1$]
\label{prz:ga9-fs-integral}
Dla $n=1$ oraz $z=x+iy$ mamy
$\omega_{\rm FS}=\frac1\pi\frac{\mathrm dx\wedge\mathrm dy}
{(1+|z|^2)^2}$. Punkt w nieskończoności nie wnosi
dodatkowego składnika do całki po zwartej sferze,
bo jest zbiorem miary zero. Zatem we współrzędnych
biegunowych
\[
 \int_{\PP^1(\C)}\omega_{\rm FS}
 =\frac1\pi\int_0^{2\pi}\int_0^\infty
 \frac{r}{(1+r^2)^2}\,\mathrm dr\,\mathrm d\theta
 =2\int_0^\infty\frac{r}{(1+r^2)^2}\,\mathrm dr=1.
\]
W ten sposób sprawdziliśmy normalizację:
$\int_{\PP^1}\omega_{\rm FS}=\deg\mathcal O(1)=1$.
Ten rachunek wiąże rzutowy stopień algebraiczny
z całką po rzeczywistej powierzchni wymiaru dwa.
\end{przyklad}

\begin{wniosek}[Każda gładka rozmaitość rzutowa jest Kählerowska]
Jeśli $X\hookrightarrow\PP^N_\C$ jest gładką, zamkniętą
podrozmaitością algebraiczną, to $X^{\rm an}$ jest
zwartą rozmaitością Kählerowską; jej formą
Kählerowską może być restrykcja $\omega_{\rm FS}$.
\end{wniosek}
\begin{proof}
Zanurzenie algebraiczne indukuje holomorficzne
zanurzenie analityczne. Forma Kählerowska na
$\PP^N(\C)$ jest zamknięta i dodatnia na
każdej zespolonej podprzestrzeni stycznej, więc
jej restrykcja do $X^{\rm an}$ zachowuje obie własności.
Zwartość wynika z domkniętości $X^{\rm an}$
w zwartej przestrzeni $\PP^N(\C)$.
\end{proof}

\section{Kiedy rozmaitość Kählerowska jest algebraiczna?}
\label{sec:ga9-kodaira}

\begin{twierdzenie}[Twierdzenie Kodairy o zanurzeniu — wynik zewnętrzny]
\label{tw:ga9-kodaira}
Niech $M$ będzie zwartą rozmaitością zespoloną.
Wtedy $M$ dopuszcza holomorficzne zanurzenie jako
zamknięta podrozmaitość pewnej $\PP^N(\C)$ wtedy
i tylko wtedy, gdy istnieje holomorficzna wiązka
liniowa $L\to M$ z dodatnio określoną rzeczywistą
formą krzywizny $(1,1)$ pewnej metryki hermitowskiej.
Taka wiązka nazywa się \emph{dodatnią}.
\end{twierdzenie}

\begin{uwaga}[Dlaczego wynik nie jest tautologią?]
Z jednej strony restrykcja $\mathcal O_{\PP^N}(1)$
na zamkniętą podrozmaitość projektową jest dodatnia.
W drugą stronę potrzebujemy wykazać istnienie tylu
sekcji $L^{\otimes m}$, by rozdzielały wszystkie
punkty oraz kierunki styczne. To jest istotna treść
twierdzenia Kodairy, oparta na teorii analitycznej
$\bar\partial$ i odpowiednich twierdzeniach o znikaniu.
Nie wyprowadzamy jej tutaj wyłącznie z dodatniości
krzywizny, nie dowodząc tych wyników pośrednich.
\end{uwaga}

\begin{wniosek}[Kryterium całkowitej klasy Kählerowskiej]
\label{wn:ga9-integral-kahler}
Zwarta rozmaitość zespolona $M$ jest projektowa
w sensie holomorficznego zanurzenia do przestrzeni
rzutowej wtedy i tylko wtedy, gdy ma formę Kählerowską
$\omega$, której klasa de Rhama $[\omega]$ należy
do obrazu $H^2(M,\Z)\to H^2(M,\R)$.
\end{wniosek}
\begin{proof}
Jeśli $M$ jest projektowa, restrykcja znormalizowanej
formy Fubiniego--Study'ego jest Kählerowska, a jej
klasa jest całkowita, bo jest pierwszą klasą Cherna
wiązki $\mathcal O(1)|_M$.
W drugą stronę twierdzenie o wykładniczym ciągu snopów,
połączone z faktem, że rzeczywista klasa $[\omega]$
jest typu $(1,1)$, daje holomorficzną wiązkę liniową
$L$ z $c_1(L)=[\omega]$ w kohomologii rzeczywistej.
Dzięki lematowi $\partial\bar\partial$ na zwartej
rozmaitości Kählerowskiej można dobrać metrykę
hermitowską na $L$, której znormalizowana krzywizna
jest dokładnie $\omega$. Jest więc ona dodatnia,
a twierdzenie Kodairy daje zanurzenie rzutowe.
W bardziej szczegółowym rozumowaniu całkowitą klasę w\,~$H^2(M,\Z)$
(reprezentującą $[\omega]$ po przejściu do $H^2(M,\R)$)
uznajemy za klasę typu $(1,1)$. Twierdzenie Lefschetza o klasach
$(1,1)$, wyprowadzone z ciągu wykładniczego i rozkładu Hodge'a,
dostarcza wtedy $L$. Lemat $\partial\bar\partial$ zmienia
metrykę na $L$ tak, by jej forma Cherna była równa $\omega$.
Oba te kroki opierają się na twierdzeniach zewnętrznych,
które dokładniej omówimy w następnym rozdziale.
\end{proof}

\begin{przyklad}[Zwarte tory zespolone: Kähler nie oznacza rzutowości]
Niech $\Lambda\subset\C^n$ będzie kratą rzędu $2n$,
a $T=\C^n/\Lambda$ zespolonym torusem. Standardowa
forma $\omega_0$ jest niezmiennicza względem translacji,
więc przechodzi na dodatnią zamkniętą formę na $T$:
każdy torus zespolony jest Kählerowski.
Jednak dla $n\ge2$ istnieją tory zespolone bez
jakiejkolwiek całkowitej klasy Kählerowskiej,
a zatem nieprojektowe. Na bardzo ogólnym torusie
zespolonym o wymiarze dwa zachodzi nawet
$H^{1,1}(T)\cap H^2(T,\Z)=\{0\}$ (po przejściu do
kohomologii bez torsji). Istnienie takich struktur
można wykazać argumentem z przeliczalnością warunków
na macierz okresów; szczegółowy dowód tego istnienia
pozostawiamy jako wynik zewnętrzny.
\end{przyklad}

\section{Hodge, Dolbeault i niezmienniki topologiczne}
\label{sec:ga9-hodge}

\begin{twierdzenie}[Dolbeault — wynik zewnętrzny]
Niech $M$ będzie zwartą rozmaitością zespoloną.
Jeśli $\Omega^p_M$ oznacza snop holomorficznych
$p$-form, to dla każdego $p,q\ge0$ istnieje
naturalny izomorfizm
$H^q(M,\Omega^p_M)\simeq H^{p,q}_{\bar\partial}(M)$,
gdzie prawa strona jest kohomologią kompleksu
gładkich form typu $(p,\bullet)$ z operatorem $\bar\partial$.
\end{twierdzenie}

\begin{twierdzenie}[Rozkład Hodge'a dla zwartych rozmaitości Kählerowskich]
\label{tw:ga9-hodge}
Niech $M$ będzie zwartą rozmaitością Kählerowską.
Wtedy dla każdego $k\ge0$ zachodzi naturalny rozkład
\[
 H^k_{\rm dR}(M;\C)
 =\bigoplus_{p+q=k}H^{p,q}_{\bar\partial}(M),
 \qquad \overline{H^{p,q}}=H^{q,p}.
\]
W szczególności, jeśli $h^{p,q}=\dim_\C H^{p,q}_{\bar\partial}(M)$,
to $b_k=\sum_{p+q=k}h^{p,q}$ i $h^{p,q}=h^{q,p}$.
\end{twierdzenie}

\begin{uwaga}[O co trzeba uzupełnić dowód?]
Kluczowe są tożsamości Kählerowskie, z których
wynika zależność $\Delta_{\rm d}=2\Delta_{\bar\partial}$
(zgodnie ze standardową normalizacją laplasjanów).
Dzięki temu harmoniczna forma de Rhama rozkłada się
na składowe typu $(p,q)$, również harmoniczne.
Pełny dowód tożsamości Kählerowskich i teorii
harmonicznych reprezentantów korzysta z analizy
eliptycznej; w tym rozdziale przyjmujemy go
jako wynik zewnętrzny, a poniżej dowodzimy jego
konkretnych konsekwencji.
\end{uwaga}

\begin{wniosek}[Nieparzyste liczby Bettiego są parzyste]
\label{wn:ga9-odd-betti}
Jeśli $M$ jest zwartą rozmaitością Kählerowską,
to dla każdego $r\ge0$ liczba Bettiego $b_{2r+1}(M)$
jest parzysta. W szczególności $b_1(M)=2h^{1,0}(M)$.
\end{wniosek}
\begin{proof}
Z rozkładu Hodge'a
$b_{2r+1}=\sum_{p+q=2r+1}h^{p,q}$.
Ponieważ $p+q$ jest nieparzyste, nie może być $p=q$.
Składniki możemy zatem pogrupować w rozłączne pary
$(p,q)$ i $(q,p)$ o tych samych wymiarach. Każda para
wnosi liczbę parzystą. Dla $k=1$ pozostają wyłącznie
$(1,0)$ oraz $(0,1)$, co daje drugi wzór.
\end{proof}

\begin{przyklad}[Dlaczego nie wszystkie zwarte rozmaitości zespolone są Kählerowskie?]
Klasyczna powierzchnia Hopfa otrzymywana jako iloraz
$\C^2\setminus\{0\}$ przez działanie
$(z,w)\mapsto(2z,2w)$ jest zwartą rozmaitością zespoloną.
Jako rozmaitość gładka jest dyfeomorficzna z
$S^1\times S^3$, więc jej pierwsza liczba Bettiego
wynosi $b_1=1$. Na mocy
wniosku~\ref{wn:ga9-odd-betti} nie może być Kählerowska.
W szczególności nie jest projektowa. Zwartość
i zespolona struktura nie wymuszają więc algebraiczności.
Identyfikację topologiczną z $S^1\times S^3$
można zobaczyć, zapisując $\C^2\setminus\{0\}
\cong S^3\times\R$ przez promień $\|z\|$ i kierunek
$z/\|z\|$; mnożenie przez $2$ przesuwa współrzędną
$\log\|z\|$ o $\log2$.
\end{przyklad}

\begin{przyklad}[Hodge na krzywej eliptycznej]
Niech $E$ będzie gładką rzutową krzywą zespoloną
rodzaju jeden. Z rozdziału~\ref{ch:krzywe-snopy}
wiemy, że $h^0(E,\Omega_E^1)=1$, a dualność
Serre'a daje $h^1(E,\mathcal O_E)=1$.
Wymiary Hodge'a wynoszą
$h^{0,0}=h^{1,1}=1$, $h^{1,0}=h^{0,1}=1$.
Stąd $b_0=1$, $b_1=2$, $b_2=1$, a
$\chi_{\rm top}(E)=b_0-b_1+b_2=0$.
Warto porównać ten wynik z zerową charakterystyką
Eulera torusa rzeczywistego $S^1\times S^1$.
\end{przyklad}

\begin{twierdzenie}[Twarde Lefschetza — wynik zewnętrzny]
\label{tw:ga9-hard-lefschetz}
Niech $M$ będzie zwartą rozmaitością Kählerowską
o wymiarze zespolonym $n$, a $L$ operatorem
$\alpha\mapsto[\omega]\smile\alpha$ w kohomologii.
Dla każdego $0\le k\le n$ mnożenie przez
$[\omega]^{n-k}$ indukuje izomorfizm
$L^{n-k}:H^k(M;\R)\to H^{2n-k}(M;\R)$.
\end{twierdzenie}

\begin{wniosek}[Ograniczenia topologii gładkich odmian rzutowych]
Niech $X$ będzie gładką rzutową rozmaitością nad $\C$
o wymiarze zespolonym $n$. Jej nieparzyste liczby
Bettiego są parzyste, a dla $k<n$ zachodzi
$b_k(X)\le b_{k+2}(X)$. Ponadto klasy podrozmaitości
algebraicznych kodujące przecięcia są zgodne z
mnożeniem przez klasę Kählerowską.
\end{wniosek}
\begin{proof}
Pierwszy wniosek otrzymujemy z rozkładu Hodge'a.
Aby wykazać nierówność, użyjmy twierdzenia
\ref{tw:ga9-hard-lefschetz}. Jeśli $k<n$, a
$[\omega]\smile\alpha=0$ dla klasy
$\alpha\in H^k(M;\R)$, to tym bardziej
$[\omega]^{n-k}\smile\alpha=0$. Twarde Lefschetza
pokazuje, że $\alpha=0$. Mnożenie przez $[\omega]$
jest więc iniekcją $H^k\hookrightarrow H^{k+2}$.
Zgodność dla klas podrozmaitości wynika z
interpretacji $c_1(\mathcal O(1))$ jako klasy
hiperpłaszczyzny i z formuły Cherna--Weila;
jej ogólny aspekt enumeratywny jest częścią
twierdzeń porównawczych teorii przecięć i
kohomologii, przyjętych tutaj zewnętrznie.
\end{proof}

\section{Rachunek łączący stopień, całkę i rodzaj krzywej}
\label{sec:ga9-degree-genus}

\begin{stwierdzenie}[Stopień krzywej przez całkę]
Niech $C\hookrightarrow\PP^2_\C$ będzie gładką,
spójną, rzutową krzywą płaską stopnia $d\ge1$.
Wówczas
$\int_{C^{\rm an}}\omega_{\rm FS}=d$,
gdzie całkujemy restrykcję znormalizowanej formy
Fubiniego--Study'ego. Ponadto
$2g(C)-2=d(d-3)$.
\end{stwierdzenie}
\begin{proof}
Pierwszą klasą Cherna wiązki $\mathcal O_C(1)$
jest restrykcja klasy hiperpłaszczyzny $h$.
Przez utożsamienie jej z $[\omega_{\rm FS}]$
oraz definicję stopnia wiązki liniowej otrzymujemy
\[
 \int_{C^{\rm an}}\omega_{\rm FS}
 =\langle c_1(\mathcal O_C(1)),[C]\rangle
 =\deg\mathcal O_C(1)=d.
\]
W ostatnim kroku przecięcie z ogólną prostą w
$\PP^2$ ma $d$ punktów, liczonych z krotnościami,
co wynika z Bézouta.

Ze wzoru adjunkcji na gładkiej krzywej mamy
$\omega_C\cong
(\omega_{\PP^2}\otimes\mathcal O_{\PP^2}(C))|_C
\cong\mathcal O_C(d-3)$,
ponieważ $\omega_{\PP^2}\cong\mathcal O_{\PP^2}(-3)$.
Stopień wiązki kanonicznej $C$ wynosi $2g(C)-2$,
a $\deg\mathcal O_C(d-3)=d(d-3)$. Po przyrównaniu
stopni dostajemy drugi wzór.
\end{proof}

\begin{przyklad}[Trzy stopnie, trzy topologie]
Dla $d=1,2,3$ gładkie spójne krzywe płaskie mają
odpowiednio rodzaj $g=0,0,1$. W każdym przypadku
całka restrykcji $\omega_{\rm FS}$ wynosi jednak
odpowiednio $1,2,3$. Stopień zanurzenia i rodzaj
krzywej mierzą różne aspekty geometrii: pierwszy
zależy od wiązki $\mathcal O_C(1)$, drugi od
wewnętrznej kanoniki $K_C$. Z tego powodu tę samą
krzywą abstrakcyjną można zanurzać w przestrzeniach
rzutowych na różne sposoby, uzyskując inne stopnie.
\end{przyklad}

\section{Stein, Cartan i granica między afinicznością a analitycznością}
\label{sec:ga9-stein}

\begin{definicja}[Rozmaitość Steinowska]
Zespoloną rozmaitość $M$ nazywamy \emph{Steinowską},
gdy jest holomorficznie wypukła i holomorficznie
separowalna: dla każdego zwartego $K\subset M$
jego otoczka holomorficzna
$\widehat K_{\mathcal O(M)}=
 \{x\in M:|f(x)|\le\sup_{y\in K}|f(y)|\ \forall f\in\mathcal O(M)\}$
jest zwarta, a dowolne dwa różne punkty $x,y\in M$
można rozdzielić globalną funkcją holomorficzną
$f\in\mathcal O(M)$ z $f(x)\ne f(y)$.
\end{definicja}

\begin{twierdzenie}[Afiniczne odmiany nad $\C$ są Steinowskie — wynik standardowy]
\label{tw:ga9-affine-stein}
Jeśli $X$ jest afiniczną rozmaitością algebraiczną
nad $\C$, to $X^{\rm an}$ jest zespoloną
przestrzenią Steinowską (w przypadku gładkim:
rozmaitością Steinowską).
\end{twierdzenie}
\begin{proof}[Uzasadnienie przez zanurzenie do przestrzeni afinicznej]
Ponieważ $X$ jest afiniczna i skończonego typu,
można ją przedstawić jako domknięty podschemat
$\A^N_\C$. Zatem jej analityfikacja jest
zamkniętą przestrzenią analityczną w $\C^N$.
Przestrzeń $\C^N$ jest Steinowska, a zamknięte
podprzestrzenie analityczne przestrzeni Steinowskich
są Steinowskie. Ostatni krok jest standardowym
twierdzeniem analitycznym i nie wynika tylko
z samej definicji wypukłości.
\end{proof}

\begin{twierdzenie}[Cartan B — wynik zewnętrzny]
Jeżeli $X$ jest przestrzenią Steinowską,
a $\mathcal F$ koherentnym analitycznym snopem
$\mathcal O_X$-modułów, to
$H^q(X,\mathcal F)=0$ dla wszystkich $q\ge1$.
\end{twierdzenie}

\begin{uwaga}[Algebraiczna i analityczna afiniczność to nie to samo]
Dla afinicznego schematu $X$ oraz quasi-koherentnego
snopu $\mathcal F$ zachodzi algebraiczne
$H^q(X,\mathcal F)=0$ dla $q>0$. Dla $X^{\rm an}$
dostajemy analityczny odpowiednik przez Cartana B.
Zgodność dwóch twierdzeń w tej sytuacji nie oznacza
jednak GAGA: już na $\A^1$ grupy globalnych funkcji
regularnych i holomorficznych są różne.
W szczególności analityczność grafu funkcji $e^z$
nie czyni go algebraicznym podzbiorem $\A^2$.
Nie wolno też wnioskować bez dodatkowych założeń,
że każda rozmaitość algebraiczna, której analityfikacja
jest Steinowska, musi być afiniczna.
\end{uwaga}

\section{Synteza: kiedy teoria algebraiczna i analityczna się zgadzają?}
\label{sec:ga9-podsumowanie}

Najbardziej użyteczne rozróżnienie można zebrać
w czterech przypadkach. Na przestrzeni afinicznej
algebraiczne funkcje regularne są zwykle tylko częścią
funkcji holomorficznych. Na przestrzeni rzutowej
zamknięte zbiory analityczne są algebraiczne dzięki
Chow. Dla przestrzeni właściwych nad $\C$ snopy
koherentne i ich kohomologia są porównywalne dzięki
GAGA. Dla gładkich rozmaitości rzutowych geometria
Fubiniego--Study'ego udostępnia dodatkowo narzędzia
Kählera, Hodge'a i Lefschetza.

\begin{uwaga}[Co jest dowiedzione, a co przyjęte?]
W tym rozdziale od podstaw wykazaliśmy niealgebraiczność
grafu wykładniczego, lokalność pierścienia kiełków,
rachunek Čecha na $\PP^1$, dodatniość formy
Fubiniego--Study'ego, jej całkę na $\PP^1$,
parzystość nieparzystych liczb Bettiego przy
założeniu twierdzenia Hodge'a oraz związki
stopnia krzywej z całkami i kanoniką.
Chow został wyprowadzony z jawnie wskazanych
silniejszych wyników: GAGA, koherencji Oki--Cartana
i twierdzeń Serre'a. Zewnętrzne pozostają m.in.:
przygotowanie Weierstrassa, GAGA, Cartan B,
Cherna--Weil, tożsamości Kählerowskie,
Kodaira, Hodge i twarde Lefschetza.
\end{uwaga}

% Literatura podstawowa i źródła kontroli:
% J.-P. Serre, Géométrie algébrique et géométrie analytique (GAGA), 1956.
% P. Griffiths, J. Harris, Principles of Algebraic Geometry.
% C. Voisin, Hodge Theory and Complex Algebraic Geometry I.
% R. O. Wells, Differential Analysis on Complex Manifolds.
% The Stacks Project, Complex Analytic Spaces and GAGA:
% https://stacks.math.columbia.edu/


% ROZDZIAL 10. GEOMETRIA ALGEBRAICZNA A TOPOLOGIA
% Zaleznosci: rozdzialy 2, 4, 6, 7 oraz 9.
